\begin{figure}[ht]
    \centering
    \begin{tikzpicture}
        \begin{axis}[
            width=8.2cm, height=7cm,
            xlabel={$10^{-3} J \, [kg \, m^2]$},
            ylabel={Gęstość prawdopodobieństwa},
            grid=major,
            scaled x ticks=false,
            x tick label style={/pgf/number format/fixed, /pgf/number format/precision=4},
            legend style={at={(1.00,1.00)}, anchor=north east, font=\footnotesize}
        ]

        % --- Zaleznosc I ---
        \def\muOne{2.97}
        \def\sigmaOne{0.44}

        % --- Zaleznosc II ---
        \def\muTwo{2.25}
        \def\sigmaTwo{0.34}

        % --- Zaleznosc III ---
        \def\muThree{2.38}
        \def\sigmaThree{0.35}

        % --- Szacowanie ---
        \def\muFour{4.23}
        \def\sigmaFour{0.49}

        \def\medium{3}
        \def\range{1.5}

        % --- Zaleznosc I ---
        \addplot [
            domain=\medium - \range : \medium + \range,
            samples=1000,
            color=red,
            thick,
            smooth
        ]
        {1/(\sigmaOne*sqrt(2*pi)) * exp(-((x-\muOne)^2)/(2*\sigmaOne^2))};
        \addlegendentry{Zależność I}

        % --- Zaleznosc II ---
        \addplot [
            domain=\medium - \range : \medium + \range,
            samples=1000,
            color=blue,
            thick,
            smooth
        ]
        {1/(\sigmaTwo*sqrt(2*pi)) * exp(-((x-\muTwo)^2)/(2*\sigmaTwo^2))};
        \addlegendentry{Zależność II}

        % --- Zaleznosc III ---
        \addplot [
            domain=\medium - \range : \medium + \range,
            samples=1000,
            color=green,
            thick,
            smooth
        ]
        {1/(\sigmaThree*sqrt(2*pi)) * exp(-((x-\muThree)^2)/(2*\sigmaThree^2))};
        \addlegendentry{Zależność III}

        % --- Szacowanie ---
        \addplot [
            domain=\medium - \range : \medium + \range,
            samples=1000,
            color=purple,
            thick,
            smooth
        ]
        {1/(\sigmaFour*sqrt(2*pi)) * exp(-((x-\muFour)^2)/(2*\sigmaFour^2))};
        \addlegendentry{Szacowanie}

        % --- Punkty --- %
        \pgfmathsetmacro{\peakOne}{1/(\sigmaOne*sqrt(2*pi))}
        \draw [dashed, red!70] (axis cs:\muOne, 0) -- (axis cs:\muOne, \peakOne);
        \addplot[mark=*, mark size=1.5pt, red] coordinates {(\muOne, \peakOne)};
        \addplot[mark=*, mark size=1.5pt, red] coordinates {(\muOne, 0)};
        
        \pgfmathsetmacro{\peakTwo}{1/(\sigmaTwo*sqrt(2*pi))}
        \draw [dashed, blue!70] (axis cs:\muTwo, 0) -- (axis cs:\muTwo, \peakTwo);
        \addplot[mark=*, mark size=1.5pt, blue] coordinates {(\muTwo, \peakTwo)};
        \addplot[mark=*, mark size=1.5pt, blue] coordinates {(\muTwo, 0)};

        \pgfmathsetmacro{\peakThree}{1/(\sigmaThree*sqrt(2*pi))}
        \draw [dashed, green!70] (axis cs:\muThree, 0) -- (axis cs:\muThree, \peakThree);
        \addplot[mark=*, mark size=1.5pt, green] coordinates {(\muThree, \peakThree)};
        \addplot[mark=*, mark size=1.5pt, green] coordinates {(\muThree, 0)};
        
        \pgfmathsetmacro{\peakFour}{1/(\sigmaFour*sqrt(2*pi))}
        \draw [dashed, purple!70] (axis cs:\muFour, 0) -- (axis cs:\muFour, \peakFour);
        \addplot[mark=*, mark size=1.5pt, purple] coordinates {(\muFour, \peakFour)};
        \addplot[mark=*, mark size=1.5pt, purple] coordinates {(\muFour, 0)};

        \end{axis}
    \end{tikzpicture}
    \caption{Porównanie gęstości prawdopodobieństwa uzyskanych momentów bezwładności dla każdej z trzech metod wraz z zaznaczonymi wartościami oczekiwanymi.}
    \label{img:gestosc-prawdop-porownanie}
\end{figure}
